% this file is called up by thesis.tex
% content in this file will be fed into the main document

\chapter{Conclusion}\label{ch:chapter6}


% ----------------------- paths to graphics ------------------------

% change according to folder and file names
\ifpdf
    \graphicspath{{6_conclusions/figures/PNG/}{6_conclusions/figures/PDF/}{6_conclusions/figures/}}
\else
    \graphicspath{{6_conclusions/figures/EPS/}{6_conclusions/figures/}}
\fi


% ----------------------- contents from here ------------------------


\section{Summary}
Models x, y, z
% ---------------------------------------------------------------------------
\section{Conclusions}
% DRAFT
One of the more striking findings of this study concerns the gap between the variety of data available and the narrowness of what turned out to matter. The pipeline drew on features spanning every major domain of ageing, yet prediction converged on 20 indicators from physical health and functional status alone. This points to something worth reflecting on: self-rated health, pain, and medication use are not just convenient proxies for disease risk. In survey cohorts like \gls{TILDA}, they are integrative measures, absorbing the accumulated effect of social circumstances, lifestyle, psychological state, and biological vulnerability into a single reported experience of burden. Whether that signal can be disaggregated into clinically actionable components is a question that requires richer data than what the publicly available release provides.
% use the below numbers to streghten the abpve paragrpah using stats etc..
% in chapter 1 we say: assess whether domain specific features contribute independent predictive value beyond the core health burden set. so its imporant right to touch on this exmaple look at how many thematic domain variables exist from our chapter3 - domain_counts_heatmap.png - down below is latex representation of this figure.
% \begin{table}
% \caption{Number of features per thematic domain across waves}
% \label{tab:domain_counts}
% \begin{tabular}{lrrrrrrrr}
% \toprule
%  & SocioDemographic\_Context & PhysicalHealth\_Clinical & FunctionalStatus\_Frailty & Lifestyle\_Behavioural & Nutrition\_Metabolic & Psychological\_Cognitive & Other\_Administrative & Total\_Features \\
% \midrule
% wave1 & 272 & 149 & 48 & 58 & 10 & 97 & 175 & 809 \\
% wave2 & 119 & 148 & 18 & 21 & 0 & 38 & 138 & 482 \\
% wave3 & 164 & 184 & 44 & 14 & 8 & 93 & 132 & 639 \\
% wave4 & 161 & 186 & 30 & 13 & 4 & 36 & 162 & 592 \\
% wave5 & 128 & 186 & 34 & 14 & 0 & 27 & 202 & 591 \\
% Total\_per\_Domain & 844 & 853 & 174 & 120 & 22 & 291 & 809 & 3113 \\
% \bottomrule
% \end{tabular}
% \end{table}

\subsection{Data availability Statement} 
The data analysed in this study is subject to the following licenses/restrictions: The datasets generated during and/or analysed during the current study are not publicly available due to data protection regulations but are accessible at TILDA on reasonable request. 
Requests to access these datasets should be directed to https://tilda.tcd.ie/data/accessing-data/.  
% use some of the info from this: This analysis was conducted using the publicly accessible \gls{TILDA} dataset (Waves~1--5). The public release is a heavily pseudonymised dataset covering self-reported health, physical, cognitive, nutritional, and behavioural variables collected at each wave, together with \texttt{ICD10} coded disease indicators available from Wave~2 onwards. Several data sources relevant to chronic disease prediction are not included in the public release and require either on site hotdesk access at \gls{TILDA}'s Trinity College Dublin facility, or remote access through the \gls{TILDA}~VISTA trusted research environment. These include a dedicated biomarkers dataset, the Wave~3 \gls{MRI} dataset, an epigenetics dataset, linked primary care records (\gls{TILDA}-\texttt{PCRS}), and linked mortality records (\gls{TILDA}~Mortality). Access to these sources would enable future work to distinguish more precisely between perceived health burden, objective physiological deterioration, and clinically verified disease trajectories. The inclusion of family history indicators in Track~B partly reflects the constraints of the public release: individual-level disease diagnoses specific to the participant were limited in the available data, and family reported diagnoses served as an indirect proxy for certain conditions where participant level onset data was absent or sparse.

\subsection{Ethics Statement}  
The studies involving human participants were reviewed and approved by the \texttt{Faculty of Health Sciences Research Ethics Committee} at Trinity College Dublin, Ireland. 
The patients/participants provided their written informed consent to participate in this study.

% ---------------------------------------------------------------------------
\section{Contributions}

% ---------------------------------------------------------------------------
\section{Future Work}
To use wave6 published on 15/12/2025, IDS and covid-19 datasets. 
How well our models generalise well compared to international data.
The utilisation of the harmonised dataset.



% ---------------------------------------------------------------------------
% ----------------------- end of thesis sub-document ------------------------
% ---------------------------------------------------------------------------